\chapter{总结与展望}\label{chap:conclusion}

\section{总结}

总结部分应回顾研究的主要工作和成果，概括本文的核心贡献。建议按照研究目标逐条归纳主要结论，简明扼要地阐述每项工作的意义。例如：

\begin{enumerate}
    \item 本文研究了……问题，提出了……方法，并通过实验验证了其有效性。
    \item 在……方面，本文改进了……，相比已有方法，在……指标上提升了……。
    \item 本文设计并实现了……系统/算法/模型，实验结果表明……。
\end{enumerate}

以上条目仅作为结构参考，撰写时应根据实际研究内容调整。

\section{展望}

展望部分应针对研究中尚未解决的问题，提出后续研究的方向和建议。可以讨论现有方法的局限性，以及可能的改进思路和应用前景。例如：

\begin{enumerate}
    \item 本文方法在……场景下存在一定局限性，未来可通过……加以改进。
    \item 随着……技术的发展，本文的研究成果可进一步应用于……领域。
    \item 后续工作可考虑引入……方法，以提升……方面的性能。
\end{enumerate}
